\subsection{Spin transport and the information needed to select exchange statistics}
\label{sec:spin-exchange-selection-boundary}

A one-particle spin representation and a many-particle exchange rule are
different data. Likewise, a current obtained from a supplied fermionic
action need not distinguish that action from a bosonic alternative on the
preparations actually used. The following finite statements separate these
two questions without assuming a relativistic field reconstruction.

\begin{proposition}[A spin lift does not select a two-particle exchange sector]
\label{prop:spin-exchange-sector-boundary}
Let $\mathcal H$ be a finite-dimensional complex Hilbert space with
$\dim\mathcal H\geq2$, and let $D:G\to U(\mathcal H)$ be a supplied
unitary representation. Suppose a central element $z\in G$ acts as
$D(z)=-I$. On $\mathcal H\otimes\mathcal H$, let
$\tau(v\otimes w)=w\otimes v$. Both eigenspaces
\[
 \mathcal H_+=\ker(\tau-I)=\operatorname{Sym}^2\mathcal H,
 \qquad
 \mathcal H_-=\ker(\tau+I)=\bigwedge^2\mathcal H
\]
are nonzero invariant subspaces of $D\otimes D$. The central element acts
as $+I$ on both. For every nonzero $v\in\mathcal H$, the vector
$v\otimes v$ is a nonzero member of $\mathcal H_+$, whereas its
antisymmetric projection is zero. Thus the one-particle spin action,
including its central minus sign, does not by itself determine the
exchange sign or exclude two particles in the same one-particle mode.
\end{proposition}
\begin{proof}
For any linear operator $A$ on $\mathcal H$,
\[
 \tau(A\otimes A)=(A\otimes A)\tau,
\]
as is checked on simple tensors, which span the tensor product. Hence
the orthogonal projections $\Pi_\pm=(I\pm\tau)/2$ commute with
$D(g)\otimes D(g)$ for every $g$. The restrictions are therefore unitary
representations. Since $(-I)\otimes(-I)=I$, they have the stated central
action. If $v\ne0$, then $\|v\otimes v\|=\|v\|^2>0$ and
$\tau(v\otimes v)=v\otimes v$, so $\Pi_-(v\otimes v)=0$.
For orthonormal $u,v$, the unit vector
$(u\otimes v-v\otimes u)/\sqrt2$ belongs to $\mathcal H_-$.
Both sectors therefore exist for the same supplied spin representation,
but they disagree on the doubled-mode state.
\end{proof}

This statement applies, in particular, to a supplied two-dimensional
unitary realization of a binary rotation group whose central element is
$-I$. A nonsplit double cover or a unique spin structure does not remove
the displayed choice of two-particle sector. The statement does not
construct a relativistic bosonic spinor field satisfying locality and a
positive-energy spectrum condition, and it is not a counterexample to a
relativistic spin--statistics theorem. It establishes only the insufficiency
of the finite one-particle representation data just specified.

\begin{proposition}[Conserved single-occupancy channels cannot select statistics]
\label{prop:single-channel-statistics-boundary}
Let
\[
 \mathcal K=\bigoplus_{\lambda\in\Lambda}\mathcal K_\lambda
\]
be a finite orthogonal decomposition into one-particle channel spaces.
In both the bosonic and fermionic Fock spaces over $\mathcal K$, retain
the sector with exactly one particle in every channel. Use the same
initial normalized channel vectors $\psi_\lambda$, the same classical
registers, and the same ordered finite sequence of the following operations:
\begin{enumerate}
 \item evolution by a self-adjoint channel-preserving quadratic Hamiltonian,
 or second quantization of
 a specified channel-preserving one-particle unitary;
 \item readout of expectations of channel-preserving one-body operators;
 \item classical register updates whose values are determined by those
 readouts and the retained classical registers.
\end{enumerate}
The two statistics give identical one-body densities, all such readouts,
and every resulting classical register value. In particular, an
expectation current, mean Gauss residual, and classical feedback history
formed from these data cannot determine the exchange statistics.
\end{proposition}
\begin{proof}
Order the finite channel set once. On the algebraic tensor product define
\[
 J_\pm\Bigl(\bigotimes_\lambda v_\lambda\Bigr)
   =\prod_\lambda a_\pm^\dagger(v_\lambda)\,|0\rangle,
\]
where $+$ denotes the bosonic and $-$ the fermionic creation operators,
and the product follows the fixed order. Orthogonality of the channel
spaces and the canonical relations show that the inner product of two
such vectors is $\prod_\lambda\langle v_\lambda,w_\lambda\rangle$
in both cases. The maps therefore extend to unitary identifications
between $\bigotimes_\lambda\mathcal K_\lambda$ and the respective
fixed-channel-occupation sectors.

For a one-body operator $B_\lambda$ supported in channel $\lambda$,
its number-conserving second quantization obeys
\[
 J_\pm^*\,d\Gamma_\pm(B_\lambda)\,J_\pm
  = I\otimes\cdots\otimes B_\lambda\otimes\cdots\otimes I.
\]
In the fermionic calculation, the signs from moving annihilation and
creation through the preceding occupied channels cancel. In the bosonic
calculation there is no sign, and occupation one gives the same matrix
coefficient. Consequently the restriction of a channel-diagonal
quadratic Hamiltonian $d\Gamma_\pm(\bigoplus_\lambda h_\lambda)$
is the same sum of tensor-slot Hamiltonians. The restrictions of
second-quantized channel-preserving unitaries are likewise the same
tensor products. Starting from the same tensor-product vector therefore
gives the same channel vectors after every such operation. Their common
one-body density is
\[
 P=\bigoplus_\lambda
       |\psi_\lambda\rangle\langle\psi_\lambda|,
 \qquad
 \langle d\Gamma_\pm(B)\rangle=\operatorname{Tr}(PB)
\]
for channel-diagonal $B$. Equal classical inputs select equal subsequent
operations and equal readouts give equal classical updates. Induction on
the operation sequence proves the assertion, including adaptive
mean-field feedback. A channel declared to remain empty contributes its
vacuum factor and zero one-body density in both constructions.
\end{proof}

For the finite hypercharge-current construction, take
$\lambda=(s,a)$ with $a=1,\ldots,d_s$ and
$\mathcal K_\lambda=\mathbb C^{64}\otimes\mathbb C^2$, the prepared
site and Spin coordinates. There is one occupied orbital in each of the
fifteen channels. Its abelian hopping and phase interventions act within
each channel. Replacing the supplied CAR Fock state by the bosonic number
state with one particle in each of the same channels preserves
\[
 P=\bigoplus_s I_{d_s}\otimes
       |\psi_s\rangle\langle\psi_s|,
 \qquad 0\leq P\leq I,\qquad \operatorname{Tr}P=15.
\]
The projector property in this preparation is therefore not itself a
statistics discriminator. The multiplicity index is an orthogonal
one-particle label; repeating the same spatial--Spin orbital across
different multiplicity indices does not prepare two particles in one
complete one-particle mode.

The local Cayley updates use the same one-particle matrices in the two
constructions. This comparison second-quantizes each one-particle Cayley
unitary; it does not replace it by the Cayley transform of the full
many-body Hamiltonian. The midpoint orbitals, and hence the midpoint
variational functional and its current, also agree algebraically. Those
midpoint averages are not normalized intermediate Fock states. Therefore
the charge transfers, electric kicks, and expectation Gauss residuals
coincide. The subsequent electric drifts and Wilson plaquette kicks use
only the common classical registers, so they preserve this equality
before the next matter factor reads the evolved link.

The retained site-charge variance supplies no distinction either. A
single particle in channel $(s,a)$ has site probability
$p_{s,v}=\|\psi_{s,v}\|^2$, and its site-number projector squares to
itself for both statistics. The product over channels gives
\[
 \operatorname{Var}(Q_v)
   =\sum_s d_s q_s^2p_{s,v}(1-p_{s,v})
\]
in both states. The nonzero variance excludes an operator Gauss
constraint with the specified classical electric numbers. Its equality
does not turn that constraint into a consequence of mean Gauss.

These identical histories concern the executed abelian, channel-preserving
restriction. They do not assert equivalence of the full nonabelian or
Yukawa theories, of channel-mixing evolutions, or of preparations with two
particles in one conserved channel. Nor do they change the validity of
the CAR current calculation under its stated assumptions. They show why
that current calculation alone does not select those assumptions. An
exchange-sensitive comparison must specify additional preparation and
readout operations outside this restriction, together with their source
realization. A finite comparison of supplied alternatives is distinct
from deriving physical spin--statistics from observer consistency.

\begin{proposition}[A two-particle port discriminator]
\label{prop:two-particle-port-discriminator}
Let $L,R$ be orthonormal modes of one declared channel, let
$a=3/5$, $b=4/5$, and supply the one-particle unitary
\[
 U=\begin{pmatrix}a&-ib\\-ib&a\end{pmatrix}.
\]
Prepare $\Psi_\pm=(L\otimes R\pm R\otimes L)/\sqrt2$ and apply
$U\otimes U$. For the coincidence effect
\[
 F=|LR\rangle\langle LR|+|RL\rangle\langle RL|,
\]
the antisymmetric and symmetric preparations have probabilities
\[
 p_-=1,\qquad p_+=(a^2-b^2)^2=\frac{49}{625}.
\]
A distinguishable-particle preparation $L\otimes R$ gives
$p_{\rm dist}=a^4+b^4=337/625$. All three preparations have the same
summed one-particle density $I$ before and after the operation, so this
contrast requires a two-particle readout.
\end{proposition}
\begin{proof}
The matrix is unitary because $a^2+b^2=1$, and its determinant is one.
Direct expansion gives
\[
 (U\otimes U)\Psi_-=\Psi_-,\qquad
 (U\otimes U)\Psi_+
  =(a^2-b^2)\Psi_+-i\sqrt2 ab\,(LL+RR).
\]
The three vectors $\Psi_+,LL,RR$ are orthonormal, so the bosonic
coincidence probability is $(a^2-b^2)^2$, with probability
$2a^2b^2=288/625$ at each bunched port. The antisymmetric output has no
bunched component. For the distinguishable product input the coincidence
amplitudes are $a^2$ and $-b^2$ in orthogonal labelled outcomes; their
squared moduli add to $337/625$. Its bunching probability at each port is
$a^2b^2=144/625$. Finally the two orthogonal occupied orbitals give summed
one-particle density $I$, and $UIU^\dagger=I$.
\end{proof}

The port labels denote a supplied invariant channel basis. They do not
assert that the two physical ports carry identical raw spin-coordinate
vectors: a spin-dependent hopping matrix can use different local basis
vectors. The rational coefficients agree with a finite hopping factor,
but this discriminator supplies its own two-particle preparation and
coincidence effect. It inherits no electric flux dressing, operator Gauss
constraint, clock identification or laboratory instrument from a
one-particle-per-channel history.

Indistinguishability is essential. If the particles carry normalized
unobserved internal vectors $\alpha,\beta$, unchanged by the port unitary,
write $\eta=|\langle\alpha,\beta\rangle|^2$. Expanding the normalized
symmetrized or antisymmetrized input and tracing over those vectors gives
\[
 p_{\rm fermion}=\frac{337+288\eta}{625},\qquad
 p_{\rm boson}=\frac{337-288\eta}{625}.
\]
The cross term is the product
$\langle\alpha,\beta\rangle\langle\beta,\alpha\rangle$; its sign is
the exchange sign. Orthogonal hidden states give $\eta=0$ and erase the
contrast. The ideal values above require $\eta=1$.

There is also a conditional stability bound. Let $\rho$ be either ideal
output density and $\widetilde\rho$ an actual normalized output density
with $\|\widetilde\rho-\rho\|_1/2\leq\epsilon$. If an actual effect
$0\leq\widetilde F\leq I$ satisfies
$\|\widetilde F-F\|\leq\delta$, then
\[
 \left|\operatorname{Tr}(\widetilde\rho\widetilde F)
       -\operatorname{Tr}(\rho F)\right|
 \leq\epsilon+\delta.
\]
Indeed, split the difference into
$\operatorname{Tr}((\widetilde\rho-\rho)F)$ and
$\operatorname{Tr}(\widetilde\rho(\widetilde F-F))$.
The positive and negative parts of the trace-zero density difference
each have trace at most $\epsilon$, so the first term has magnitude at
most $\epsilon$; the second is bounded by $\delta$. A preparation
trace-distance bound transfers unchanged through an exactly common
unitary; imperfect dynamics instead require a bound on the actual output.
For a common error budget $\epsilon+\delta<288/625$, the two ideal-model
probability intervals are disjoint because their separation is $576/625$.
These are supplied error bounds, not a calibration or a bound on finite
sampling error. An observed test additionally needs authenticated
preparation, indistinguishability, operation and readout evidence.
